%%%%%%%%%%%%%%%%%%%%%%%%%%%%%%%%%%%%%%%%%%%%%%%%%%%%%%%%%%%%%%%%%%%%%%%%%%%%%%%
% hka_symbols.tex -- zentrale Geometrie der HKA-Symbole
%
% Die Pfade fuer H, K und A stehen ausschliesslich hier.
% Layout-Dateien bestimmen nur Farbe, Position und Skalierung.
%%%%%%%%%%%%%%%%%%%%%%%%%%%%%%%%%%%%%%%%%%%%%%%%%%%%%%%%%%%%%%%%%%%%%%%%%%%%%%%

% H -------------------------------------------------------------------------
% #1 Farbe, #2 TikZ-Ankerkoordinate, #3 Skalierung
\NewDocumentCommand{\HKASymbolH}{O{ThemeColor} m O{0.45}}{%
	\begin{scope}[shift={(#2)}]
		\begin{scope}[scale=#3]
			\path [fill=#1, rounded corners=1]
				(0,-2.95) -- (3,-2.95) -- (3,0) -- (4.2,0) -- (4.2,-2.95)
				-- (7.2,-2.95) -- (7.2,-3.95) -- (4.2,-3.95) -- (4.2,-6.9)
				-- (3,-6.9) -- (3,-3.95) -- (0,-3.95) -- cycle;
			\path [fill=#1, rounded corners=1]
				(7.8,0) -- (9,0) -- (9,-6.9) -- (7.8,-6.9) -- cycle;
		\end{scope}
	\end{scope}%
}

% K --------------------------------------------------------------------------
% #1 Farbe, #2 TikZ-Ankerkoordinate, #3 Skalierung
\NewDocumentCommand{\HKASymbolK}{O{ThemeColor} m O{0.45}}{%
	\begin{scope}[shift={(#2)}]
		\begin{scope}[scale=#3]
			\path [fill=#1, rounded corners=1]
				(3,0) -- (4.2,0) -- (4.2,-2.8) -- (9,0) -- (9,-1.3)
				-- (5.2,-3.45) -- (9,-5.6) -- (9,-6.9) -- (4.2,-4.1)
				-- (4.2,-6.9) -- (3,-6.9) -- cycle;
		\end{scope}
	\end{scope}%
}

% A --------------------------------------------------------------------------
% #1 Farbe, #2 TikZ-Ankerkoordinate, #3 Skalierung
\NewDocumentCommand{\HKASymbolA}{O{ThemeColor} m O{0.45}}{%
	\begin{scope}[shift={(#2)}]
		\begin{scope}[scale=#3]
			\path [fill=#1, rounded corners=1]
				(6,0) -- (6.65,0) -- (10.05,-6.9) -- (8.75,-6.9) -- (6,-1.3)
				-- (4.8,-3.7) -- (7.52,-3.7) -- (8.85,-4.8) -- (4.3,-4.8)
				-- (3.25,-6.9) -- (1.95,-6.9) -- (5.35,0) -- cycle;
		\end{scope}
	\end{scope}%
}

% Kompositionen ---------------------------------------------------------------
% Diese Befehle enthalten keine eigene Geometrie, sondern kombinieren nur die
% drei obigen Grundformen.
\NewDocumentCommand{\HKASymbolKA}{O{ThemeColor} m m O{0.45}}{%
	\HKASymbolK[#1]{#2}[#4]%
	\HKASymbolA[#1]{#3}[#4]%
}

\NewDocumentCommand{\HKASymbolFull}{O{ThemeColor} m m m O{0.45}}{%
	\HKASymbolH[#1]{#2}[#5]%
	\HKASymbolK[#1]{#3}[#5]%
	\HKASymbolA[#1]{#4}[#5]%
}
